\chapter{Application Pratique : Réalisation de l'Audit LiZbiZ}
\label{chap:appaudit}

Ce chapitre met en œuvre les concepts théoriques sur le terrain. Nous simulons la phase de réalisation de l'audit (Section 6.4 de l'ISO 19011) au sein de l'entreprise LiZbiZ.

\section{Étape 1 : Collecte des Preuves}

L'auditeur doit collecter des informations vérifiables pour étayer ses conclusions.

\subsection{Les Trois Sources de Preuves}

\begin{enumerate}
    \item \textbf{Entretiens (Interviews) :} Échanges avec le personnel. (Subjectif, à recouper).
    \item \textbf{Observation :} Vue directe des activités. (Objectif, instant T).
    \item \textbf{Examen Documentaire et Technique :} Analyse de procédures, de logs, de configurations. (Objectif, historique).
\end{enumerate}

\begin{boxdef}
\textbf{Piste d'audit (Audit Trail) :} Enchaînement logique et documenté des preuves permettant de reconstituer l'historique d'un événement ou d'une transaction.
\end{boxdef}

\subsection{Technique de l'Échantillonnage}

Il est impossible de vérifier l'intégralité des logs ou des 150 comptes utilisateurs. L'auditeur pratique l'\textbf{échantillonnage}.

\begin{itemize}
    \item \textbf{Échantillonnage statistique :} Choix aléatoire de 10 comptes.
    \item \textbf{Échantillonnage orienté risque :} Choix ciblé (ex: les comptes admins, les nouveaux arrivés, les comptes inactifs depuis longtemps).
\end{itemize}

% ------------------------------------------------------------------------------
\section{Étape 2 : Audit Technique sur le Lab GNS3}
% ------------------------------------------------------------------------------

Nous connectons maintenant l'auditeur au laboratoire GNS3 de LiZbiZ pour un audit technique réel.

\subsection{Scénario 1 : Audit de la Segmentation Réseau}

\textbf{Critère d'audit (ISO 27001 - A.5.15) :} Le réseau doit être segmenté pour isoler les flux sensibles.

\begin{boxlizbiz}
\textbf{Mission Pratique :}
L'auditeur se connecte au routeur/firewall (pFsense) via la console.

\textbf{Tâche :} Vérifier les règles de filtrage inter-VLAN.
\textbf{Commande / Action :} Visualiser les règles (Access Control Lists).

\textbf{Constat attendu :}
L'auditeur découvre une règle "Allow Any -> Any" entre le VLAN 40 (IoT Logistique) et le VLAN 20 (Applicatif).
\textbf{Preuve :} Capture d'écran de la règle.
\textbf{Analyse :} Si un scanner IoT est compromis, il peut atteindre l'ERP. Non-respect du principe de moindre privilège.
\end{boxlizbiz}

\subsection{Scénario 2 : Audit des Droits d'Accès}

\textbf{Critère d'audit (ISO 27001 - A.5.18) :} Les droits d'accès doivent être révoqués lors du départ d'un employé.

\begin{boxlizbiz}
\textbf{Mission Pratique :}
L'auditeur accède au contrôleur de domaine (Windows Server AD).

\textbf{Tâche :} Lister les utilisateurs actifs et comparer avec la liste des employés fournie par la RH (document fictif "ListeEmployes\_Janvier.pdf").

\textbf{Constat attendu :}
L'auditeur trouve le compte "jmartin" actif. La liste RH indique que "Jean Martin" a quitté l'entreprise il y a 3 mois.
\textbf{Preuve :} Capture d'écran de l'Active Directory montrant "Enabled = True".
\textbf{Analyse :} Procédure de départ non appliquée. Risque de compte orphelin.
\end{boxlizbiz}

% ------------------------------------------------------------------------------
\section{Étape 3 : Audit Organisationnel (Jeu de Rôle)}
% ------------------------------------------------------------------------------

L'audit n'est pas que technique. Il faut auditer les processus humains.

\subsection{Scénario : Audit du Processus de Gestion des Incidents}

\textbf{Critère d'audit (ISO 27001 - A.5.24) :} L'organisation doit gérer les incidents de sécurité.

\begin{boxlizbiz}
\textbf{Mise en situation :}
L'auditeur interroge le Responsable Helpdesk (joué par un étudiant).

\textbf{Question de l'Auditeur :} "Que s'est-il passé le 12 janvier dernier suite à l'alerte antivirus sur le poste de la comptabilité ? Pouvez-vous me montrer le ticket ?"

\textbf{Réponse simulée (Erreur) :} "Ah, on a juste supprimé le virus. C'était un faux positif je pense, on n'a pas jugé utile d'ouvrir un ticket pour ça."

\textbf{Analyse de l'Auditeur :}
Absence de traceabilité. La procédure exige un ticket pour tout incident.
\end{boxlizbiz}

% ------------------------------------------------------------------------------
\section{Étape 4 : Analyse des Écarts et Formulation des Constats}
% ------------------------------------------------------------------------------

Une fois les preuves collectées, il faut formuler les constats.

\subsection{La Structure d'un Constat d'Audit}

Un bon constat répond à la structure : \textbf{Critère - Constat - Preuve - Cause - Risque}. Il ne doit jamais être une opinion.

\begin{table}[htbp]
\centering
\small
\begin{tabular}{p{3cm}p{9cm}}
    \toprule
    \textbf{Élément} & \textbf{Exemple LiZbiZ (Cas compte orphelin)} \\
    \midrule
    \textbf{Critère} & ISO 27001 A.5.18 : "Les droits d'accès doivent être retirés lors du départ". \\
    \textbf{Constat} & Le compte "jmartin" est toujours actif 3 mois après le départ. \\
    \textbf{Preuve} & Capture AD + Liste RH. \\
    \textbf{Cause} & Absence de communication fluide entre RH et IT. \\
    \textbf{Risque} & Utilisation malveillante des données par un ancien employé. \\
    \bottomrule
\end{tabular}
\caption{Formulation d'un constat d'audit}
\end{table}

\subsection{Classification de la Non-Conformité}

L'auditeur doit qualifier la gravité de l'écart.

\begin{itemize}
    \item \textbf{Non-Conformité Majeure :} Absence totale de système ou défaillance généralisée. (Ex: Aucune procédure de gestion des accès n'existe).
    \item \textbf{Non-Conformité Mineure :} Incident isolé, ponctuel. (Ex: Le compte "jmartin" est le seul oublié, la procédure existe mais a eu un raté).
    \item \textbf{Observation :} Piste d'amélioration, pas de non-respect strict. (Ex: La procédure de départ est sur papier, il serait mieux de la dématérialiser).
\end{itemize}

\begin{boxlizbiz}
\textbf{Décision pour LiZbiZ :}
L'oubli du compte "jmartin" est classé \textbf{Non-Conformité Mineure}, car la procédure existe. Cependant, si l'auditeur trouve 10 comptes orphelins sur 150, cela deviendra une \textbf{Non-Conformité Majeure} (défaillance du système de contrôle).
\end{boxlizbiz}

% ==============================================================================
% Fichier : 03_module_audit.tex (Extrait Chapitre 5)
% Description : Rapport d'Audit et Plan d'Actions Correctives
% ==============================================================================